\documentclass[11pt]{article}
\usepackage[T1]{fontenc}
\usepackage[margin=1in]{geometry}
\usepackage{amsmath,amssymb,amsthm,mathtools}
\usepackage{graphicx}
\usepackage{microtype}
\usepackage[hidelinks]{hyperref}

\newtheorem{theorem}{Theorem}
\newtheorem{lemma}[theorem]{Lemma}
\newtheorem{proposition}[theorem]{Proposition}
\newtheorem{remark}[theorem]{Remark}
\newcommand{\R}{\mathbb R}
\newcommand{\D}{\mathcal D}
\newcommand{\E}{\mathcal E}
\newcommand{\EE}{\mathbb E}
\newcommand{\eps}{\varepsilon}

\title{Non-Gaussian Tilts and Gaussian Rigidity in\\Stereographic Infinite-Volume Limits}
\author{Candidate full resolution for arXiv:1701.05569}
\date{13 August 2026}

\begin{document}
\maketitle

\begin{center}
\fbox{\parbox{0.91\linewidth}{\textbf{Status.}
Candidate full resolution, likely valid; human review required.  The
unrestricted compactness construction contains explicit non-Gaussian limits,
but every translation-invariant limit under the source's uniform $L^2$
hypothesis is the free massive Gaussian.  Novelty is provisional.}}
\end{center}

\section{Source problem and precise answer}

Zahariev constructs infinite-volume probability measures on
$\D'(\R^d)$ by transporting normalized density perturbations of Gaussian
measures on the sphere through stereographic projection and scaling.  The
abstract of \cite{source} asks whether this construction contains
non-Gaussian measures:

\begin{figure}[ht]
\centering
\includegraphics[width=\linewidth]{figures/open_problem_crop.png}
\caption{The exact open problem in the abstract of \cite{source}.}
\end{figure}

There are two natural readings because Proposition~4.1 of the source allows
arbitrary densities satisfying norm bounds, whereas its Euclidean-invariant
Glimm--Jaffe output imposes symmetry and reflection-positive structure.  We
settle both:

\begin{enumerate}
\item In the literal unrestricted class, the answer is \emph{yes}: a bounded
cosine tilt gives an explicit non-Gaussian limit.
\item In the advertised translation-invariant/Glimm--Jaffe class, the answer
is \emph{no}: uniform $L^2$ contiguity and ergodicity force the limit to be the
free massive Gaussian.
\end{enumerate}

\section{Setup from the source}

Let $I_k=U_\alpha U_{\beta_k}:\D(\R^d)\to\D(\mathbb S^d)$ be the source's
test-function transport and let
\[
 R_k:\D'(\mathbb S^d)\longrightarrow\D'(\R^d),\qquad
 (R_k\phi)(f)=\phi(I_kf),
\]
be its transpose.  Write $\gamma_k^S$ for the centered sphere Gaussian with
smoothed covariance $\widetilde C_{S,k}$, and suppose
\begin{equation}\label{eq:bounds}
 \inf_k\|\rho_k\|_{L^1(\gamma_k^S)}\ge K_1>0,
 \qquad
 \sup_k\|\rho_k\|_{L^2(\gamma_k^S)}\le K_2<\infty.
\end{equation}
The normalized interacting and transported measures are
\[
 \mu_k^S=N_k\rho_k\gamma_k^S,
 \qquad N_k^{-1}=\int\rho_k\,d\gamma_k^S,
 \qquad
 \nu_k=(R_k)_\#\mu_k^S,
 \quad \bar\gamma_k=(R_k)_\#\gamma_k^S.
\]

The exact covariance identity and mollifier estimate in the free-field remark
of \cite{source} imply
\begin{equation}\label{eq:freeconv}
 \bar\gamma_k\Longrightarrow\gamma,
 \qquad
 \widehat\gamma(f)=
 \exp\!\left[-\frac12\left\langle
 f,(-\Delta+m^2)^{-1}f\right\rangle_{L^2}\right],
\end{equation}
where $m>0$.  Proposition~4.1 gives subsequential limits of $\nu_k$.

\section{Proof intuition}

There is an easy way to preserve a local non-Gaussian perturbation: make the
sphere density depend on one field coordinate that is transported to a fixed
Euclidean test function.  This proves the literal positive answer, but it
chooses a location and breaks Euclidean invariance.

The symmetry requirement changes everything.  The source's Cauchy--Schwarz
estimate says that each interacting measure is uniformly $L^2$-dominated by
its Gaussian reference.  Because both sequences have been transported to one
space, this estimate survives weak convergence and produces an $L^2$ density
of the limit with respect to the free Gaussian.  The latter is mixing under
translations.  A translation-invariant probability with such a density can
only have the constant density one.

\section{The unrestricted positive example}

\begin{proposition}\label{prop:tilt}
The hypotheses \eqref{eq:bounds} admit a sequence whose transported measures
converge to a non-Gaussian probability on $\D'(\R^d)$.
\end{proposition}

\begin{proof}
Fix $0<\eps<1$ and a real nonzero $h\in\D(\R^d)$.  Define on sphere fields
\[
 \rho_k(\phi)=1+\eps\cos\bigl(\phi(I_kh)\bigr).
\]
This is a continuous positive density and
$1-\eps\le\rho_k\le1+\eps$, so \eqref{eq:bounds} holds.  Since
$\phi(I_kh)=(R_k\phi)(h)$, for every bounded continuous $F$ on
$\D'(\R^d)$,
\[
 \int F\,d\nu_k
 =\frac{\displaystyle\int F(\Phi)
       \bigl(1+\eps\cos\Phi(h)\bigr)\,d\bar\gamma_k(\Phi)}
       {\displaystyle\int
       \bigl(1+\eps\cos\Phi(h)\bigr)\,d\bar\gamma_k(\Phi)}.
\]
Both numerator and denominator converge by \eqref{eq:freeconv}; the
denominator is at least $1-\eps$.  Hence
\begin{equation}\label{eq:tiltlimit}
 \nu_k\Longrightarrow\nu_h,
 \qquad
 \frac{d\nu_h}{d\gamma}(\Phi)
 =\frac{1+\eps\cos\Phi(h)}
 {\int(1+\eps\cos\Psi(h))\,d\gamma(\Psi)}.
\end{equation}
The Gaussian variable $X=\Phi(h)$ has positive variance.  Its law under
$\nu_h$ is the normal law multiplied by $1+\eps\cos x$ and renormalized.
This is not Gaussian: the Radon--Nikodym ratio between two equivalent
one-dimensional Gaussian laws is an exponential quadratic, whereas the
positive bounded oscillatory function $1+\eps\cos x$ is not.  Thus $\nu_h$
is non-Gaussian.
\end{proof}

\begin{remark}
The density in Proposition~\ref{prop:tilt} is deliberately non-invariant and
does not claim to satisfy the source's special reflection-positivity
hypotheses.  It answers exactly the unrestricted Proposition~4.1 reading.
\end{remark}

\section{Uniform domination survives the limit}

\begin{lemma}\label{lem:domination}
Suppose along a subsequence that $\nu_k\Rightarrow\nu$ and
$\bar\gamma_k\Rightarrow\gamma$.  Under \eqref{eq:bounds}, there is
$r\in L^2(\gamma)$ such that
\[
 r\ge0,\qquad \int r\,d\gamma=1,\qquad \nu=r\gamma,
 \qquad \|r\|_2\le K_2/K_1.
\]
\end{lemma}

\begin{proof}
The Cauchy--Schwarz estimate used in the appendix of \cite{source} gives, for
every bounded continuous $F$ on $\D'(\R^d)$,
\begin{equation}\label{eq:domk}
 \int|F|\,d\nu_k
 \le \frac{K_2}{K_1}
 \left(\int|F|^2\,d\bar\gamma_k\right)^{1/2}.
\end{equation}
Indeed, apply the source estimate to $F\circ R_k$ before transport.  Weak
convergence passes \eqref{eq:domk} to the limit for bounded continuous $F$.

Let $\mathcal C$ be the linear span of bounded continuous cylinder functions.
The functional $L(F)=\int F\,d\nu$ on $\mathcal C$ is therefore bounded in
$L^2(\gamma)$ with norm at most $K_2/K_1$.  Cylinder functions are dense in
$L^2(\gamma)$, so $L$ extends uniquely.  Riesz representation gives
$r\in L^2(\gamma)$ with $L(F)=\int Fr\,d\gamma$.  Positivity of $L$ gives
$r\ge0$, and $L(1)=1$ gives $\int r\,d\gamma=1$.  Taking the cylinder
functions $F(\Phi)=e^{i\Phi(f)}$ shows that $\nu$ and $r\gamma$ have the same
characteristic functional.  Probability measures on $\D'(\R^d)$ are
determined by that functional, hence $\nu=r\gamma$.
\end{proof}

\section{Mixing of the massive free field}

Let $\tau_a$ translate distributions by $a\in\R^d$.  The convention is
irrelevant as long as $\Phi(f)\circ\tau_a=\Phi(f_a)$ for the oppositely
translated test function.

\begin{lemma}\label{lem:mixing}
The Gaussian probability $\gamma$ in \eqref{eq:freeconv} is strongly mixing
under translations.  In particular, its translation-fixed subspace in
$L^2(\gamma)$ consists of the constants.
\end{lemma}

\begin{proof}
For $f,g\in\D(\R^d)$, the cross-covariance after a translation is, up to the
Fourier normalization,
\[
 \left\langle f,(-\Delta+m^2)^{-1}g_a\right\rangle
 =\int_{\R^d}e^{ip\cdot a}
 \frac{\overline{\widehat f(p)}\widehat g(p)}{|p|^2+m^2}\,dp.
\]
The integrand without the phase is in $L^1$, so this tends to zero as
$|a|\to\infty$ by Riemann--Lebesgue.  The Gaussian characteristic-function
formula then shows that correlations of cylinder exponentials
$e^{i\Phi(f)}$ and their translates factor in the limit.  Their linear span
is dense in $L^2(\gamma)$; approximation proves strong mixing for all
$L^2$ functions.

If $q\in L^2(\gamma)$ is translation fixed, mixing gives
$\|q\|_2^2=\langle q,q\circ\tau_a\rangle\to|\int q\,d\gamma|^2$.
Thus $q$ has zero variance and is constant.
\end{proof}

\section{Gaussian rigidity and the full answer}

\begin{theorem}\label{thm:main}
Let $\nu$ be any subsequential scaling limit in the construction of
\cite{source} under \eqref{eq:bounds}.  If $\nu$ is translation invariant,
then
\[
 \nu=\gamma,
\]
the centered massive free Gaussian measure with covariance
$(-\Delta+m^2)^{-1}$.
\end{theorem}

\begin{proof}
Lemma~\ref{lem:domination} gives $\nu=r\gamma$ with $r\in L^2(\gamma)$.
Both $\nu$ and $\gamma$ are translation invariant.  Radon--Nikodym uniqueness
therefore gives $r\circ\tau_a=r$ in $L^2(\gamma)$ for every $a\in\R^d$.
Lemma~\ref{lem:mixing} makes $r$ constant, and $\int r\,d\gamma=1$ makes that
constant equal to one.
\end{proof}

The source proves that $O(d+1)$-invariant sphere densities yield
Euclidean-invariant limits.  Its bounded continuous self-interactions have
this symmetry.  Theorem~\ref{thm:main} therefore says that every such output
under the exact uniform $L^2$ hypotheses is the free field.  In particular,
the source's conjectured recovery of a non-Gaussian $P(\phi)_2$ plane measure
cannot occur without leaving those exact hypotheses or changing the reference
scheme.

\section{Verification, scope, and novelty}

The argument does not assume that absolute continuity is automatically
closed under weak limits; Lemma~\ref{lem:domination} uses the uniform $L^2$
estimate that makes it closed here.  Translation invariance is essential, as
Proposition~\ref{prop:tilt} demonstrates.

Bounded exact-title, exact-question, citation, non-Gaussianity, and
stereographic-scaling searches through 13 August 2026 found no explicit later
resolution of this scheme.  Duch--Dybalski--Jahandideh \cite{ddj} construct a
non-Gaussian $P(\Phi)_2$ limit from growing spheres via stochastic
quantization and different moment estimates.  Dybalski \cite{dybalski}
computes nontrivial connected functions after a different field
renormalization.  Neither supplies the uniform-$L^2$ classification above.
Novelty should nevertheless be regarded as provisional pending specialist
review.

\section{Human-review recommendation}

Review this as a full resolution.  The highest-value checks are: (i) the
transpose transport in Proposition~\ref{prop:tilt}; (ii) passage of
\eqref{eq:domk} to the limit and the cylinder-density argument; and (iii) the
scope boundary between Proposition~4.1 and the source's Euclidean-invariant
Glimm--Jaffe theorem.

\begin{thebibliography}{9}
\raggedright
\enlargethispage{4\baselineskip}

\bibitem{source}
S. Zahariev,
\emph{Infinite volume limits in Euclidean quantum field theory via
stereographic projection},
arXiv:1701.05569 (2017; revised 2018).

\bibitem{ddj}
P. Duch, W. Dybalski, and A. Jahandideh,
\emph{Stochastic quantization of two-dimensional $P(\Phi)$ Quantum Field
Theory},
arXiv:2311.04137 (2023).

\bibitem{dybalski}
W. Dybalski,
\emph{Interacting Euclidean quantum fields from the Gaussian free field},
arXiv:2502.07546 (2025).

\end{thebibliography}

\end{document}
